\section{Background}
\label{sec:background}

The design of accountable, autonomous multi-agent systems demands a precise answer to a foundational question: when one agent delegates a task to another, through what structural mechanism is the chain of accountability preserved? This question becomes acute in systems that support \emph{recursive spawning} — the capacity for an agent to decompose its assigned task by provisioning child agents that operate with delegated authority. As agentic runtimes migrate from single-agent to hierarchical multi-agent architectures \cite{wooldridge1995,stone2000}, the problem of managing the lifecycle of spawned agents — and of maintaining an unbroken accountability chain from every active agent back to an originating human decision — transitions from an architectural detail to a core safety requirement.

This section situates the Recursive Spawning Protocol (RSP) within five distinct research and engineering traditions: agent lifecycle management in multi-agent systems; accountability and provenance in distributed autonomous systems; the structural distinction between fixed and mutable delegation chains; hierarchical task decomposition in classical and modern AI; and the BX3 Framework's three-layer model as an architectural substrate for immutable parent chains. Together, these traditions define the conceptual landscape RSP inhabits and the specific gaps in prior work that RSP is designed to close.

% -------------------------------------------------------
\subsection{Agent Lifecycle Management in Multi-Agent Systems}
% -------------------------------------------------------

Multi-agent systems (MAS) research has long recognized that the lifecycle of an agent — its creation, activation, operation, suspension, and termination — is a first-class design concern, not an implementation detail. Stone and Veloso's foundational taxonomy classifies multi-agent systems along axes of agent heterogeneity, knowledge quality, environment accessibility, and interaction nature (cooperative versus competitive) \cite{stone2000}. Within this taxonomy, recursive spawning is most relevant to cooperative, homogeneous systems operating in partially accessible environments — precisely the regime in which modern agentic runtimes such as AutoGen \cite{autogen2023}, LangChain \cite{langchain2023}, and CrewAI \cite{crewai2024} operate.

Ferber and Müller formalized agent lifecycle management as interactions between an agent's lifecycle state machine and its organizational context, arguing that organizational roles — not just individual agent states — must be modeled to ensure system-wide coherence \cite{ferber1996}. Padgham and Thiel demonstrated that incomplete lifecycle management, particularly the failure to distinguish between suspension and termination states, introduces systemic failures: orphaned or zombie agents continue to consume resources and emit effects after their parent has exited \cite{padgham2001}. These findings establish lifecycle management as an ongoing architectural discipline requiring coherent authority boundaries at every state transition throughout the spawn chain.

Recent work on LLM-based multi-agent systems has surfaced the specific problem of \emph{unbounded spawning}. Zhuge et al. document the ``agent chain problem'': LLM-based agents that recursively spawn children to handle sub-tasks, without a defined termination condition or depth ceiling, produce chains of dozens of intermediate agents whose aggregate authority may differ substantially from what any single agent was authorized to exercise \cite{zhuge2024}. This problem is not adequately addressed by existing frameworks, which treat agent creation as a runtime policy choice rather than a first-class authorization event with structural accountability requirements.

Modern frameworks are shifting toward treating delegation as a first-class primitive. AGENTHIVE enables decentralized control through first-class delegation: any agent may spawn and coordinate sub-agents, dynamically forming task-adaptive topologies driven by task demands rather than predefined graphs \cite{agenthive2025}. AgentSpawn introduces adaptive collaboration through dynamic spawning, where agents are created mid-task based on runtime complexity metrics, enabling context-preserving child agent creation and skill inheritance across the delegation chain \cite{agentspawn2026}. The Recursive Delegation Engine (RDE) formalizes lifecycle decisions as a utility-optimization problem: at each task node, an agent decides to \emph{execute}, \emph{decompose}, or \emph{delegate} based on learned policies over success/failure counters that propagate upward through the chain \cite{rde2024}. Agent Contracts formalize resource and temporal constraints on spawned agents — start conditions, runtime budgets (tokens, API calls, time), output contracts, and explicit termination criteria — preventing resource leaks and enabling orderly handoffs in delegation chains \cite{agentcontracts2025}.

The critical gap across this landscape: none of these frameworks treats the spawn event as a first-class accountability transfer with structurally immutable provenance. Spawning is modeled as a creation event or a policy decision; it is not modeled as a delegation event that creates a new accountable entity with an immutable chain back to a human principal.

% -------------------------------------------------------
\subsection{Accountability and Provenance in Distributed Systems}
% -------------------------------------------------------

Accountability in distributed agent systems requires two interrelated capabilities: (1) a tamper-evident record of which agent performed which action, and (2) a cryptographically verifiable chain linking each action back to its originating authorization. Without both, audit trails degenerate into siloed logs that cannot sustain regulatory scrutiny or forensic reconstruction \cite{vap2026,trusttrack2025}.

Several recent protocols provide cryptographic provenance primitives for agentic systems. Chain of Consciousness (CoC) establishes an append-only SHA-256 hash chain of lifecycle events, anchored to Bitcoin via OpenTimestamps, with a W3C Decentralized Identifier (DID) for persistent agent identity \cite{chainofconsciousness2026}. Each session boundary is bridged by a forward-commitment hash, enabling continuity proofs that link discontinuous sessions into a single verifiable identity arc. The Human Delegation Provenance (HDP) protocol defines a token-based scheme that cryptographically captures human authorization within multi-hop agent chains \cite{hdp2026}: each delegation step is appended as a signed hop in an append-only chain; verification is offline and relies solely on the issuer's Ed25519 public key and the current session identifier — no central registry or third-party trust anchors required. HDP's append-only hop architecture is the closest prior art to RSP's Immutable Parent Pointer, differing in that HDP addresses human-to-agent authorization while RSP addresses agent-to-agent recursive spawning under bounded Purpose Layer authority.

The Verifiable AI Provenance (VAP) framework and TrustTrack both establish architectural requirements for evidentiary-grade decision trails in autonomous AI systems \cite{vap2026,trusttrack2025}. Agent Execution Records (AER) propose first-class provenance primitives where each agent action is recorded as a first-class artifact with explicit causality chains \cite{aer2025}. The Warrant Certificate Authority (WCA) model extends this with auditable data provenance for AI-agent tool-call chains \cite{wca2025}. The TIBET protocol specifies transaction and interaction evidence trails for AI-to-AI and AI-to-Human interactions \cite{tibet2026}. The MAIF consortium's artifact-centric paradigm enforces trust and provenance at the level of individual agent artifacts rather than agent identities alone \cite{maif2025}.

What unites these approaches is the recognition that accountability requires a structurally immutable chain of custody. Where they differ from RSP is in the mechanism of immutability: cryptographic approaches (hashing, signing) provide authenticity guarantees for the chain's contents, but do not by themselves prevent a runtime state modification that rewrites the chain after an action occurs. RSP's contribution is to make the parent pointer immutable not by cryptographic signature but by architectural enforcement at the Fact Layer — the pointer cannot be modified because the deterministic enforcement layer does not permit the modification, not because a signature would fail to detect it.

% -------------------------------------------------------
\subsection{Fixed Versus Mutable Delegation Chains}
% -------------------------------------------------------

Delegation chains in multi-agent systems take two structurally different forms, with profoundly different accountability properties. In a \emph{fixed delegation chain}, the parent-child relationship is established at spawn time and is immutable thereafter: the child's parent pointer cannot be changed, the child's authority cannot be re-delegated except through the defined spawn protocol, and the chain terminates at a predetermined anchor (typically a human principal). In a \emph{mutable delegation chain}, the parent-child relationship is treated as a runtime state that can be updated, redirected, or reassigned as conditions change.

Mutable delegation chains offer greater operational flexibility: a child agent can be reparented to a different parent if the organizational structure changes, if a parent fails, or if load balancing requires reassignment. This flexibility, however, comes at a significant accountability cost. If the parent-child relationship is mutable at runtime, any single action by a child agent may have been authorized by one parent at one moment and a different parent at another — and the effective authorization for any given action may not be the authorization that was recorded at spawn time. The chain of provenance becomes conditional on the state of a mutable mapping, rather than fixed by an immutable record. This introduces the \emph{delegatee ambiguity problem}: after an event, it may be impossible to determine which principal effectively authorized the action, because the delegation chain has been rewritten since the action occurred.

Fixed delegation chains sacrifice this flexibility for a stronger accountability guarantee. Each child's authority is traceable to a single, immutable parent pointer, and that pointer chains unbroken back to a human anchor. There is no ambiguity about effective authorization: the warrant that activated the child encodes the complete delegation chain as it existed at the moment of activation, and that chain cannot be altered retroactively. The BX3 Framework's recursive spawning architecture exemplifies fixed delegation: each spawned Worksheet contains a hard-coded pointer to the parent's Purpose Layer, and that pointer is architecturally indestructible — not policed by a runtime check, but made structurally immutable by the Fact Layer's deterministic enforcement \cite{bx3framework,bx3framework2026}. Every descendant node can trace its accountability lineage to a human anchor regardless of spawning depth.

The tradeoff is that reparenting in a fixed delegation chain requires formal re-spawn: the child must be deprovisioned and a new spawn event initiated with the new parent. This is a deliberate design choice that prioritizes accountability over operational flexibility. RSP adopts fixed delegation as its structural default and provides the spawn-reparent mechanism as the sanctioned path for cases where reparenting is genuinely required.

% -------------------------------------------------------
\subsection{Hierarchical Task Decomposition in AI}
% -------------------------------------------------------

Hierarchical task decomposition is the mechanism by which a complex goal is segmented into a chain or tree of simpler subgoals, each assigned to a specialized agent or processing unit. It is the structural foundation of any recursive spawning-based multi-agent architecture and has deep roots in both classical AI planning and modern LLM-based agent systems.

In classical AI, Hierarchical Task Network (HTN) planning formalizes decomposition as a recursive process: a compound task is decomposed into sub-tasks, each of which is either primitive (executable directly) or compound (requiring further decomposition) \cite{htn1993,htn2000}. The decomposition is constrained by method selection rules that encode domain-specific knowledge about how to achieve goals under different conditions. The key property of HTN decomposition for accountability analysis is that it is \emph{prescriptive}: the decomposition tree is part of the plan, not a runtime observation. Each node in the HTN tree represents a decision about how to organize work, and this decision is part of the planning record.

Modern LLM-based multi-agent systems have adopted and extended this model. ReAcTree introduces hierarchical agent trees with two node types: control flow nodes that organize execution order, and agent nodes that reason, act, and expand into further subgoals \cite{reactree2025}. Dynamic tree expansion allows agent nodes to autonomously generate and refine subgoals at runtime, substantially outperforming flat sequential baselines (63\% vs. 24\% goal success on WAH-NL using Qwen2.5 72B). StructuredAgent combines online hierarchical planning using dynamic AND/OR trees with a structured memory module that tracks candidate solutions and constraints \cite{structuredagent2025}. DeepAgent uses a Hierarchical Task DAG (HTDAG) to dynamically decompose high-level objectives into interdependent sub-tasks across multiple layers, with a recursive two-stage planner-executor enabling continuous task refinement and adaptation \cite{deepagent2025}. HiPlan introduces a two-level hierarchy: global milestone action guides derived from expert demonstrations provide high-level structure, while local step-wise hints adapt past trajectory segments to current observations, correcting deviations and aligning actions with milestone objectives \cite{hiplan2025}. ReCAP addresses context drift in sequential prompting and cross-level continuity loss in hierarchical prompting through recursive context-aware reasoning \cite{recap2025}.

What connects classical HTN and modern LLM-based decomposition is the spawn tree: a hierarchical structure in which each node represents a subtask and edges represent the decomposition relationship. In both traditions, the tree is typically mutable and implicit — decomposition decisions are not recorded as first-class authorization artifacts, and child nodes are not distinguished from parent nodes in terms of their accountability status. Classical HTN subtasks are subroutine calls without independent accountability chains; LLM sub-agents are independent processes whose authorization chain is implicit in the context window rather than explicit in an immutable warrant structure.

Recursive Spawning adopts the hierarchical decomposition paradigm from both traditions but adds the critical accountability layer absent from both: each node in the decomposition tree is a first-class accountability entity with an immutable parent pointer, a defined termination condition, and a cryptographically verifiable chain back to a human anchor. Decomposition is not merely a planning technique; it is an authorization event that creates a new accountable entity. The decomposition tree is therefore not just a planning artifact but an organizational chart of delegated authority.

% -------------------------------------------------------
\subsection{The BX3 Framework as Substrate for Immutable Parent Chains}
% -------------------------------------------------------

The BX3 Framework provides the architectural substrate within which RSP operates. As defined in \cite{bx3framework,bx3framework2026}, the BX3 Framework organizes any autonomous node into three immutable functional layers — the \textit{Purpose Layer}, the \textit{Bounds Engine}, and the \textit{Fact Layer} — each defined by the properties it must maintain rather than by the type of actor occupying it. The Purpose Layer carries intent, judgment, and final human accountability; the Bounds Engine carries constrained execution within a defined operating range; the Fact Layer carries deterministic enforcement and forensic auditability. The upstream accountability guarantee of the framework states that when any node encounters a state it cannot resolve within its bounds, accountability escalates recursively upward through the system hierarchy, bypassing all machine actors, until it reaches a human anchor \cite{bx3framework2026}.

The Fact Layer plays the central role in RSP. It is the Fact Layer that creates the immutable parent pointer, issues the spawn warrant, and maintains the forensic ledger that records every spawn event. The Fact Layer's determinism is what makes the parent pointer immutable: once written, the pointer cannot be updated because the Fact Layer does not permit modifications to its authorization records. This immutability is not a policy choice; it is an architectural property of the deterministic enforcement layer.

The BX3 Framework's treatment of parent chain immutability directly enables the recursive spawning pattern. Because each spawned agent inherits a permanent, verifiable parent reference, recursion does not degrade the accountability graph. Each generation of spawned agents adds a verifiable layer to the chain without altering the integrity of prior layers. Every spawn event is a provenance-constrained operation, not merely a creation event. This distinguishes BX3's recursive spawning from ad hoc agent spawn models: the parent pointer is established at provisioning, stored in the child's sealed memory envelope, and is never mutable after the atomic Fact Layer write confirms the spawn event \cite{bx3framework2026}.

The three-layer model provides the structural basis for the upstream accountability guarantee in the spawn context specifically. When a child agent is activated via a spawn warrant, its authority is bounded by the operating range defined by the parent node's Bounds Engine. If the child encounters a condition outside its provisioned scope, it triggers the Bailout Protocol \cite{bailoutprotocol2026} — the upward exception propagation mechanism of the BX3 Framework — which bypasses all intermediate machine actors and routes the exception to the nearest human accountability anchor. The spawn chain is therefore not only an authorization structure but an escalation structure: the same chain that carries delegated authority downward carries exceptions upward when the Bounds Engine cannot resolve them, and the chain along which accountability escalates is the same immutable chain that authorized the spawn.

This coupling of authorization and escalation is the key architectural insight of Recursive Spawning. In conventional systems, authorization and escalation are separate concerns: a spawning decision is made by the parent agent, and an escalation decision is made independently when an exception occurs. In the BX3 Framework, both decisions are governed by the same immutable chain structure. There is no separate escalation infrastructure; the spawn chain is the escalation chain, by design. The result is an accountability architecture in which the question — \emph{who is responsible for this agent's actions?} — is always answerable by traversing the immutable parent chain back to a human anchor, and in which every traversal of that chain is a structurally enforced property of the runtime rather than a forensic reconstruction performed after the fact.

The convergence of BX3's Layer model with independently developed protocols (HDP, CoC, SentinelAgent's Delegation Chain Calculus) suggests that three-layer functional decomposition is not an arbitrary design choice but a natural structural pattern that emerges from the requirements of accountable, auditable, and non-deterministic-resilient multi-agent systems \cite{hdp2026,chainofconsciousness2026,sentinelagent2025}. Recursive spawning, implemented on this substrate, provides the strongest available guarantee of upstream accountability in hierarchical agent architectures.